\section{先行研究との境界と新規性の表現}
本稿の構成要素を、既知の内容、既知結果から直接従う内容、本稿で導く結果に分ける。
\begin{center}
\begin{tabular}{p{0.48\linewidth}p{0.42\linewidth}}
\toprule
要素&本稿での分類\\
\midrule
ジャイロスコピック重力メモリー&既知\cite{SerajOblak2023}\\
ソフト／スピン成分とハード／ヘリシティ成分の分解&既知\cite{SerajOblak2023,FayeSeraj2025}\\
$N_{AB}\widetilde C^{AB}$ 型ハード汎関数&既知\cite{SerajOblak2023,FayeSeraj2025}\\
ヘリシティ／局所双対性生成子としての解釈&既知\cite{SerajOblak2023}\\
Jacobi に基づく持続的観測量／非線形平面波&既知\cite{Flanagan2019,Flanagan2020}\\
平面波メモリーテンソル&既知\cite{Harte2025}\\
Bondi--Sachs における曲線偏差観測量の階層&既知\cite{GrantNichols2022}\\
ホロノミーから導くジャイロスコピック・メモリー&既知\cite{SerajNeogi2023}\\
Bel--Robinson 曲率代数&既知\cite{FerrandoSaez2012}\\
面積制約付き曲線最適化・曲げエネルギー&隣接する既知の変分問題\cite{Hahmann2005,Ferone2016}\\
偏光面積とハード・ジャイロ汎関数の係数対応&既知汎関数からの規約固定した導出\\
無拘束 Jacobi 最小曝露則 $1/\|K_0\|$&本稿で導出\\
定常 Jacobi 最小作用則 $1/\|P_0K_0P_0\|$&本稿で導出\\
固定ヌル窓付きハード・メモリー問題の最良係数 $2/\|P_0K_0P_0\|$&本稿で導出\\
\bottomrule
\end{tabular}
\end{center}

Flanagan--Grant--Harte--Nichols は、一般の持続的観測量と非線形平面波における Jacobi 構造を確立している\cite{Flanagan2019,Flanagan2020}。Grant--Nichols は、漸近平坦時空における曲線偏差観測量を Bondi--Sachs 枠組みで展開した\cite{GrantNichols2022}。Harte--Mieling--Oancea--Steininger は、平面波散乱を少数のメモリーテンソルで整理している\cite{Harte2025}。Seraj--Neogi は、漸近ホロノミーの Lorentz 部分とジャイロスコピック・メモリーの関係を示した\cite{SerajNeogi2023}。これらを本稿独自の構造として数えない。

数学側にも近縁問題がある。Hahmann--Sauvage--Bonneau は、閉平面曲線の面積を保つ multiresolution deformation と、その最適化で扱う曲げエネルギーを研究している\cite{Hahmann2005}。Ferone--Kawohl--Nitsch は、固定面積を持つ単純閉曲線について幾何学的な elastica 最小化を研究している\cite{Ferone2016}。これらは本稿の符号付き面積汎関数と二階微分コストに近いが、閉曲線／弧長幾何、許容境界条件、固定ヌル窓の Jacobi/Bondi 対応、平均ゼロ圧縮作用素という本稿の同条件問題とは異なる。従って、本稿は「面積制約付き曲げエネルギー最適化」一般を新規とは主張しない。

本稿は歴史的優先権を主張しない。また、双線形形式に対する Cauchy--Schwarz 不等式、コンパクト作用素ノルムの飽和、面積制約付き曲線最適化という一般的機構それ自体は既知であり、そこを新規性とは数えない。最終定理の同時条件について、入出力定常、初期シアーなし、固定した放射横断ヌル窓、Bel--Robinson 曲率作用、ハード・ジャイロ／ヘリシティ目標、最良係数を同時に扱う先行定理は文献検索で特定できなかった。本稿が新規性として主張する範囲は、物理的な入出力定常条件から平均ゼロ制約と $K_{\rm stat}=P_0K_0P_0$ を特定し、固定した補助ヌルプローブ上の曲率作用とハード・ジャイロ汎関数を同じ放射データ上の制約付き逆最適化問題として結び、exact Jacobi 回復まで含めて最良係数を確定する組合せに限定する。この文献検索結果は歴史的な世界初を保証するものではない。

\section{考察}
本稿の主結果は、確立されたハード・ジャイロ汎関数と、同じ放射データを横切る固定ヌル窓上の曲率作用の先頭摂動係数との間に最良下界を与える。ハード・ジャイロスコピック・メモリーの二次係数を指定すると、それより小さな $Q_\gamma^{(2)}$ で同じ先頭次数のハード／ヘリシティ・チャネルを生成することはできず、さらに定常条件を満たす滑らかな許容回復族がその境界へ任意に近づく。これは既知のヘリシティ汎関数の単なる記号置換ではなく、波形空間上の制約付き逆最適化を伴う。ただし、最適化の最終段階そのものは標準的な作用素ノルム問題であり、結果の学術的価値はその物理的な問題設定・制約作用素・対応関係が同条件で既知かどうかに依存する。

定常性が係数を大きく変える点も重要である。無拘束 Jacobi 最小作用問題の係数はハード・メモリーへ形式的に移せるものの、その最適化波形はバーストの入出力定常性を保証しない。物理条件を正しく課すと、許容接空間は $\mathcal U_0$ に縮み、作用素ノルムは $\|K_0\|$ から $\|P_0K_0P_0\|$ に変わる。従って最良係数は約 $21.97$ ではなく約 $71.14$ となる。この差は平均ゼロ制約による作用素圧縮そのものに由来する。

一方で射程は限定される。第一に $Q_\gamma$ は標準的な重力波エネルギーではない。第二にハード・ジャイロスコピック・メモリーは $r^{-2}$ かつ弱重力場では二次であり、現時点で検出器に直接結びつく観測量としては弱い。第三にハード／ソフト分解はフレームに依存し、本稿は初期シアーなし規約に固定される。第四に本定理は弱振幅・放射域の先頭次数に関する主張であり、有限振幅における大域的な非線形時空の最小作用関数までは扱わない。

従って本稿の結論は三層である。(i) 非線形 Jacobi メモリー・チャネルには、小目標に対する厳密で鋭い最小曝露則がある。(ii) 入出力定常を課すと、同じ構造は平均ゼロ部分空間へ圧縮した作用素 $K_{\rm stat}=P_0K_0P_0$ 上で、再び厳密で鋭い法則を持つ。(iii) 標準的な初期シアーなし Bondi 放射では、その定常 Jacobi チャネルの二次係数はハード・ジャイロスコピック・メモリーの二次係数の二倍であり、$Q_\gamma^{(2)}$ をコストとする leading-order 放射データ問題は $\kappa_{\rm H}^{\rm lead}(\psi)=2|\psi|/\|K_{\rm stat}\|$ と厳密に解ける。

\section{結論}
本稿ではヌル Jacobi 散乱に対する最小作用問題を自己完結に定式化し、まず無拘束クラスで
\[
\kappa_\Omega(\omega)=\frac{|\omega|}{\|K_0\|}+O(\omega^2)
\]
を証明した。次に、物理的な入出力定常性が潮汐制御の平均ゼロ条件を課すことを用い、制限作用素 $K_{\rm stat}=P_0K_0P_0$ に対して
\[
\kappa_{\rm stat}(\omega)=\frac{|\omega|}{\|K_{\rm stat}\|}+O(\omega^2)
\]
という鋭い法則を得た。さらに Bondi シアーの局所 STF 代数と弱い TT 波の偏光面積を通じて、$\omega_\gamma^{(2)}=-2\Delta\Psi_{\rm H}^{(2)}$ を導いた。従って、標準的な初期シアーなしフレーム、入出力定常、弱振幅かつ放射域の先頭次数では
\[
\kappa_{\rm H}^{\rm lead}(\psi)=\frac{2}{\|K_{\rm stat}\|}|\psi|,
\qquad
\frac{2}{\|K_{\rm stat}\|}=71.13876223\ldots
\]
である。ここで $\psi=\Delta\Psi_{\rm H}^{(2)}$ は二次摂動係数であり、最小化されるコストは $Q_\gamma^{(2)}$ である。この係数は leading-order TT/Bondi 放射データ問題で厳密に鋭い。一方、exact nonlinear Jacobi endpoint の値関数は
$\kappa_{\rm stat}(\omega)=|\omega|/\|K_{\rm stat}\|+O(\omega^2)$
という別の漸近定理であり、full finite-amplitude Einstein 時空におけるハード・メモリー最小作用則は本稿の主張外である。

今後の理論的課題として、固定ヌル窓という補助構造を除去できるか、有限振幅の Einstein 解空間へ値関数を拡張できるか、最終シアーにも制約を課した場合に最良係数がどう変わるかが残る。
